\subsection{定义与初步性质}

设 G 为一个群，B 为 G 的一个子群。群 \(B \times B\) 通过 \((b, b')\) . \(g = bgb'^{-1}\) 作用在 G 上，其中 \(b, b' \in B\) 且 \(g \in G\)。\(B \times B\) 在 G 上的轨道是当 \(g \in G\) 时的集合 BgB，称为关于 B 的 G 的双陪集。它们构成 G 的一个划分；相应的商记为 B\textbackslash G/B。若 C 和 C' 是双陪集，则 CC' 是双陪集的并。

定义 1. Tits 系统是一个四元组 \((\mathbf{G},\mathbf{B},\mathbf{N},\mathbf{S})\)，其中 \(\mathbf{G}\) 是一个群，\(\mathbf{B}\) 和 \(\mathbf{N}\) 是 \(\mathbf{G}\) 的两个子群，\(\mathbf{S}\) 是 \(\mathrm{N} / (\mathrm{B}\cap \mathrm{N})\) 的一个子集，并满足下列公理：

(T1) 集合 \(\mathrm{B} \cup \mathrm{N}\) 生成 \(\mathrm{G}\)，且 \(\mathrm{B} \cap \mathrm{N}\) 是 \(\mathrm{N}\) 的正规子群。

(T2) 集合 S 生成群 \(W = N/(B \cap N)\)，并且由阶为 2 的元素组成。

(T3) \(sBw \subset BwB \cup BswB\)，其中 \(s \in S\) 且 \(w \in W\)。

(T4) 对所有 \(s \in S\)，都有 \(sBs \not\subset B\)。

群 \(W = N / (B \cap N)\) 有时称为 Tits 系统 \((G, B, N, S)\) 的 Weyl 群。

备注。1) 我们将在第 5 号（定理 3 的推论）看到，如果给定 (G, B, N)，则至多存在一个 N/(B∩N) 的子集 S，使得 (G, B, N, S) 是 Tits 系统。

\begin{enumerate}
\def\labelenumi{\arabic{enumi})}
\setcounter{enumi}{1}
\tightlist
\item
  设 \((G, B, N, S)\) 是一个 Tits 系统，Z 是 G 的包含于 B 中的正规子群。令 \(G' = G/Z, B' = B/Z, N' = N/(Z \cap N)\)，并令 \(S'\) 为 S 在 \(N'/(B' \cap N')\) 中的像。立即可见，\((G', B', N', S')\) 是一个 Tits 系统。
\end{enumerate}

在本段中，设 \((\mathrm{G},\mathrm{B},\mathrm{N},\mathrm{S})\) 为一个 Tits 系统，并置 \(T=B\cap N\)、W=N/T。双陪集均指 G 关于 B 的双陪集。对于任意 \(w\in W\)，置 \(C(w)=BwB\)；这是一个双陪集。

我们将从公理 (T1) 至 (T4) 推出若干初等结论。用 \(w, w', \ldots\) 表示 W 的元素，用 \(s, s', \ldots\) 表示 S 的元素。以下关系是显然的：

\[
\mathrm{C} (1) = \mathrm{B}, \quad \mathrm{C} (w w ^ {\prime}) \subset \mathrm{C} (w). \mathrm{C} (w ^ {\prime}), \quad \mathrm{C} (w ^ {- 1}) = \mathrm{C} (w) ^ {- 1}.\tag{1}
\]

公理 (T3) 也可写成

\[
\mathrm{C} (s). \mathrm{C} (w) \subset \mathrm{C} (w) \cup \mathrm{C} (s w).\tag{2}
\]

此外，由 (1) 有 \(C(sw) \subset C(s).C(w)\)，又因 \(C(s).C(w)\) 是双陪集的并，故只有以下两种可能：

\[
\mathrm{C} (s). \mathrm{C} (w) = \left\{ \begin{array}{l l} \mathrm{C} (s w) & \text {if} \mathrm{C} (w) \not \subset \mathrm{C} (s). \mathrm{C} (w) \\ \mathrm{C} (w) \cup \mathrm{C} (s w) & \text {if} \mathrm{C} (w) \subset \mathrm{C} (s). \mathrm{C} (w). \end{array} \right.\tag{3}
\]

由 (T4)，\(\mathrm{B} \neq \mathrm{C}(s). \mathrm{C}(s)\)；在 (3) 中令 \(w = s\)，并利用关系 \(s^2 = 1\)，得到

\[
\mathrm{C} (s). \mathrm{C} (s) = \mathrm{B} \cup \mathrm{C} (s).\tag{4}
\]

此公式表明，\(B \cup C(s)\) 是 G 的一个子群。将 (4) 两边右乘 \(C(w)\)，并利用公式 (3) 以及关系

\[
\mathrm{B.C} (w) = \mathrm{C} (w),
\]

得到

\[
\mathrm{C} (s). \mathrm{C} (s). \mathrm{C} (w) = \mathrm{C} (w) \cup \mathrm{C} (s w).\tag{5}
\]

对公式 (2)、(3) 和 (5) 中出现的集合取逆，然后以 \(w^{-1}\) 代替 w，得到公式

\[
\mathrm{C} (w). \mathrm{C} (s) \subset \mathrm{C} (w) \cup \mathrm{C} (w s)\tag{(2')}
\]

\[
\mathrm{C} (w). \mathrm{C} (s) = \mathrm{C} (w s) \quad \text {   if   } \mathrm{C} (w) \not \subset \mathrm{C} (w). \mathrm{C} (s)
\]

\[
\mathrm{C} (w). \mathrm{C} (s) = \left\{ \begin{array}{l l} \mathrm{C} (w) \cup \mathrm{C} (w s) & \text {if} \mathrm{C} (w) \subset \mathrm{C} (w). \mathrm{C} (s) \end{array} \right.\tag{3'}
\]

\[
\mathrm{C} (w). \mathrm{C} (s). \mathrm{C} (s) = \mathrm{C} (w) \cup \mathrm{C} (w s).\tag{(5')}
\]

引理 1. 设 \(s_1, \ldots, s_q \in S\)，且 \(w \in W\)。我们有

\[
\mathrm{C} \left(s _ {1} \dots s _ {q}\right). \mathrm{C} (w) \subset \bigcup_ {\left(i _ {1}, \dots i _ {p}\right)} \mathrm{C} \left(s _ {i _ {1}} \dots s _ {i _ {p}} w\right),
\]

其中，\((i_1, \ldots, i_p)\) 表示区间 \([1, q]\) 中严格递增的整数序列的集合。

我们对 \(q\) 作归纳证明，\(q = 0\) 的情形是平凡的。当 \(q \geqslant 1\) 时，有 \(C(s_1 \ldots s_q).C(w) \subset C(s_1).C(s_2 \ldots s_q).C(w)\)。由归纳假设，\(C(s_2 \ldots s_q).C(w)\) 包含于各个 \(C(s_{j_1} \ldots s_{j_p}w)\) 的并中，其中

\[
2 \leqslant j _ {1} <   \dots <   j _ {p} \leqslant q.
\]

由 (T3)，集合 \(C(s_{1}).C(s_{j_{1}}\ldots s_{j_{p}}w)\) 包含于集合 \(C(s_{1}s_{j_{1}}\ldots s_{j_{p}}w)\) 与 \(C(s_{j_{1}}\ldots s_{j_{p}}w)\) 的并中。这就证明了引理。

